\section{Verification of Sufficient Conditions}
The examples below relate two assumptions in the paper to properties of the loss, the satisfaction events, and the conditional outcome distribution. A query \(r\) is a decision--context pair, and \(B_\ell\) is a uniform bound on the absolute operating loss.

\paragraph{Covering (paper Assumption~4).}
Closed balls have finite VC dimension. With finitely many outcomes, bounded functions of the outcome are vectors in a finite-dimensional cube and admit polynomial-size covers in the maximum norm. Multiplying by ball indicators preserves the covering bound because the product error is at most the sum of the two \(L_2(Q)\) errors. With finite-dimensional vector outcomes, the covering bound also holds if the requirements are affine in the outcome at each decision and the bounded loss has a fixed maximum number of affine pieces, each described by a fixed maximum number of linear inequalities. The satisfaction sets and loss subgraphs then use a bounded number of unions, intersections, and complements of balls and half-spaces, so the VC covering theorem applies \citep[Section~2.6]{vanderVaartWellner1996}. The profit constraint and loss in Section~4.3 have this structure.

\paragraph{Lipschitz smoothness (paper Assumption~2).}
For conditional densities \(f_r\) with respect to a common measure \(\mu\), the smoothness inequalities hold if \(\int |f_r-f_{r'}|\,d\mu\le K\|r-r'\|\) for a finite \(K\) and all \(r,r'\in\mathcal R\). The probability of any event changes by at most half this integral, and expected loss changes by at most \(B_\ell\) times the integral. Thus \(L_p=K/2\) and \(L_m=B_\ell K\) suffice. For finitely many outcomes the integral reduces to a sum of absolute differences between conditional probabilities.
